\chapter{问题定义与保真度门控 RSI 框架设计}\label{ch:framework}

\section{任务形式化}
\subsection{序列推荐设定}
设用户集 $\mathcal{U}$、物品集 $\mathcal{I}$，$|\mathcal{I}| =
N$。用户 $u$ 的历史交互序列为 $S_u = (i_1, i_2, \dots, i_{|S_u|})$，
按时间升序排列。序列推荐模型 $\model: S_u \mapsto \R^{N}$ 对全
部物品打分，$(\model)_i(S_u)$ 越大表示物品 $i$ 越可能成为用户的下
一次交互。训练目标为最小化真实下条目上的交叉熵：
\begin{equation}\label{eq:ce}
\mathcal{L}_{\mathrm{CE}}(\model; \Dreal) =
-\frac{1}{|\Dreal|}\sum_{(u, i) \in \Dreal}
\log \frac{\exp((\model)_i(S_u))}{\sum_{j \in \mathcal{I}}
\exp((\model)_j(S_u))},
\end{equation}
其中 $\Dreal$ 为按留一法划分的训练交互集（每用户末次交互留作测
试）。评测采用全量排序：将测试物品与全部未交互物品共同排序。

\subsection{RSI 闭环的形式化}
\begin{definition}[递归自改进闭环]\label{def:rsi}
给定初始模型 $M^{(0)}$（在 $\Dreal$ 上预训练至收敛）、生成器
$G$、校验器 $V$、筛选器 $F$ 与混合权重 $\lambda$，RSI 闭环按轮次
$r = 1, 2, \dots, R$ 迭代：
\begin{align}
\Dsyn^{(r)} &= G\big(M^{(r-1)}\big) \subseteq
\mathcal{U} \times \mathcal{I}
\label{eq:gen}\\
\tilde{\Dsyn}^{(r)} &= F\big(V(\Dsyn^{(r)}; \Dreal);\ \Dreal\big)
\label{eq:gate}\\
\Dmix &= \Dreal \cup \lambda\,\tilde{\Dsyn}^{(r)}
\label{eq:mix}\\
M^{(r)} &= \mathrm{Train}\big(M^{(r-1)}; \Dmix\big)
\label{eq:train}
\end{align}
即：上一轮模型自生成候选样本（\ref{eq:gen}），校验器与筛选器依
据真实数据 $\Dreal$ 评估并放行其中合格者（\ref{eq:gate}），
按权重 $\lambda$ 与真实数据混合（\ref{eq:mix}）后，从上一轮参数
出发继续训练得到后继模型（\ref{eq:train}）。
\end{definition}

\subsection{失败模式：朴素自训练的崩塌}
\begin{definition}[朴素自训练，naive RSI]\label{def:naive}
朴素方案取 $G$ 为每用户 top-$k$ 未交互物品、$V \equiv \top$（不
做校验）、$F$ 放行全部候选。
\end{definition}

\begin{proposition}[头部自强化]\label{prop:collapse}
设模型打分经 softmax 归一化为 $\bm{q}_u$，头部物品 $h$ 在多数用
户上的预测概率满足 $q_{u,h} \geq \delta > q_{u,t}$（$t$ 为任意长
尾物品）。则朴素方案下，每轮注入的正样本以头部为主：
$\Pr[i = h \mid (u,i) \in \Dsyn] \geq \delta^k$ 量级；训练
\eqref{eq:train} 进一步抬高头部得分，形成"生成偏头部 $\to$ 训练
强化头部 $\to$ 生成更偏头部"的正反馈。多轮迭代下
$\bm{p}_{\mathrm{syn}}$ 与 $\bm{p}_{\mathrm{real}}$ 的距离单调扩
大，最终训练分布坍缩到少数头部物品，排序退化。
\end{proposition}

该命题在第四章以实验复现（朴素方案的 NDCG@10 随轮次单调崩塌，
同时合成分布 JS 距离单调扩大）。命题亦说明：崩塌源于\emph{生成
与训练的联合动力学}，仅在筛选环节做文章而不改变生成环节的偏袒
结构，控制力有限——这正是本文同时设计生成器与门控器的原因。

\section{框架总体设计}
图~\ref{fig:framework} 给出框架全景。闭环每轮包含五个环节：
多视图一致性生成（\S\ref{sec:gen}）、双轴保真度门控
（\S\ref{sec:gate}）、加权混合（\S\ref{sec:mix}）、后继训练
（\S\ref{sec:train}）与评估反馈（\S\ref{sec:feedback}）。相对
于定义~\ref{def:rsi} 的抽象闭环，本文的关键设计在 $G$ 与
$F \circ V$ 两处。

\begin{figure}[t]
\centering
% RSI 闭环框架图 (TikZ)
\begin{tikzpicture}[
  node distance=10mm and 14mm,
  box/.style={draw, rounded corners=2pt, align=center,
              minimum height=9mm, inner sep=4pt, font=\small},
  data/.style={box, fill=blue!8},
  model/.style={box, fill=orange!15},
  gate/.style={box, fill=green!12},
  gen/.style={box, fill=red!10},
  arrow/.style={-{Stealth[length=2.2mm]}, thick},
  back/.style={-{Stealth[length=2.2mm]}, thick, dashed},
  lbl/.style={font=\scriptsize, text=black!60},
]

% 顶行: 闭环主体
\node[model] (prev) {$M^{(r-1)}$\\上一轮模型};
\node[gen, right=of prev] (views) {多视图一致性生成\\$K$ 个扰动视图 $\to$ 伪标签};
\node[gate, right=of views] (gate) {双轴保真度门控\\置信轴 + JS 预算};
\node[data, right=of gate] (mix) {加权混合\\$\Dmix = \Dreal \cup \lambda\,\tilde{\Dsyn}$};
\node[model, right=of mix] (next) {$M^{(r)}$\\后继模型};

\draw[arrow] (prev) -- node[lbl, above]{$G(\cdot)$} (views);
\draw[arrow] (views) -- node[lbl, above]{候选 $\Dsyn^{(r)}$} (gate);
\draw[arrow] (gate) -- node[lbl, above]{$\tilde{\Dsyn}^{(r)}$} (mix);
\draw[arrow] (mix) -- node[lbl, above]{$V \circ F$} (next);

% 真实数据锚定
\node[data, below=16mm of gate] (real) {$\Dreal$ 真实交互数据\\$\bm{p}_{\mathrm{real}}$, 分层 $\sigma(\cdot)$};
\draw[arrow] (real) -- node[lbl, right]{门控依据} (gate);
\draw[arrow] (real) -- node[lbl, below]{混合锚定} (mix);

% 闭环反馈
\draw[back] (next.south) -- ++(0,-7mm) -| node[lbl, pos=0.12, below]{下一轮 $r{+}1$} (prev.south);

% 评估反馈
\node[box, below=13mm of mix] (eval) {评估反馈\\NDCG/HR + 漂移监控};
\draw[back] (next.south) |- (eval.east);
\draw[back] (eval.west) -| node[lbl, pos=0.75, below]{更新预算 $\varepsilon$} (gate.south);

% 门控内部标注
\node[lbl, below=2mm of gate, text width=42mm, align=center]
  {层内 $z$ 标准化 \\ $D_{\mathrm{JS}}(\text{混合}, \text{真实}) \le \varepsilon$};

\end{tikzpicture}

\caption{保真度门控的 RSI 闭环框架。实线为数据流，虚线为反馈
路径。$M^{(r-1)}$ 在 $K$ 个扰动视图上的预测经一致性检验生成伪标
签候选；候选经置信轴（分层裕度）与分布轴（JS 预算）双重门控后
与真实数据混合，训练后继模型 $M^{(r)}$；评估反馈更新阈值与预
算。}
\label{fig:framework}
\end{figure}

\section{多视图一致性伪标签生成}\label{sec:gen}
\subsection{扰动视图}
对用户序列 $S_u$ 构造 $K$ 个随机视图：以概率 $p_{\mathrm{drop}}$
随机丢弃元素后取尾部比例 $\rho$ 的截断窗：
\begin{equation}\label{eq:view}
S_u^{(k)} = \mathrm{crop}_{\rho}\big(
\{i \in S_u : b_i^{(k)} = 0\}\big),\quad
b_i^{(k)} \sim \mathrm{Bernoulli}(p_{\mathrm{drop}}),\ k = 1..K .
\end{equation}
扰动参数 $(p_{\mathrm{drop}}, \rho)$ 取轻量档（如
$(0.1, 0.9)$），保证视图与原序列语义相近。

\subsection{一致性准入}
记 $\bm{p}^{(k)} = \mathrm{softmax}(\model(S_u^{(k)}))$（已屏蔽
历史物品）。伪标签 $(u, \hat{i})$ 的生成条件为：
\begin{equation}\label{eq:consistency}
\hat{i} = \arg\max_{j \in \mathcal{I} \setminus S_u}
p_j^{(k)}\ \ \text{对全部 } k \text{ 一致},\quad
\bar{p}_{\hat{i}} \geq \tau,\quad
\bar{p}_{\hat{i}} - \bar{p}_{(2)} \geq m,
\end{equation}
其中 $\bar{p}_{j} = \frac{1}{K}\sum_k p^{(k)}_j$，
$\bar{p}_{(2)}$ 为次高物品的平均概率。

\subsection{为何一致性优于 top-$k$ 自生成}
top-$k$ 生成直接采纳模型的头部偏好，编码的信息量为"模型已知自
己相信什么"；一致性生成要求该信念在随机扰动下保持稳定，编码的
是"模型对扰动鲁棒的信念"。后者与表示学习中的不变性原理
~\cite{tarvainen2017mean,sohn2020fixmatch,berthelot2019mixmatch}
一致：扰动一致性筛出的样本携带真实监督成分更高、头部偏袒更弱
（第四章实验对比两种生成器的漂移与最终性能）。直觉上，扰动改变
了序列上下文，若模型在不同上下文下仍给出同一预测，该预测更可能
反映用户的稳定偏好而非物品的全局流行度。

\section{双轴保真度门控}\label{sec:gate}
\subsection{置信轴：分层裕度标准化}
直接以置信度门控存在结构性偏袒：头部物品的预测置信度系统性偏
高。本文将物品按真实训练集中的流行度等频划分为 $L$ 层（层 0 为
头部五分位），在\emph{层内}对候选样本的预测裕度做 z 标准化：
\begin{equation}\label{eq:znorm}
z(u, i) = \frac{c(u, i) - \mu_{\sigma(i)}}{s_{\sigma(i)} + \epsilon},
\end{equation}
其中 $c(u,i)$ 为候选置信度（如 \eqref{eq:consistency} 的平均概
率裕度），$\mu_{\sigma(i)}, s_{\sigma(i)}$ 为物品层 $\sigma(i)$
内候选置信度的均值与标准差。准入要求 $z(u,i) \geq z_{\min}$。
层内标准化使长尾物品按\emph{其层内的相对置信度}竞争，消除头部
的结构性优势。

\subsection{分布轴：JS 距离预算}
定义自生成样本的分层分布
$\bm{p}_{\mathrm{syn}} \in \Delta^{L-1}$（各流行度层中样本占
比），真实交互的分层分布 $\bm{p}_{\mathrm{real}}$。分布轴要求
\emph{混合}分布不偏离真实分布超过预算 $\varepsilon$：
\begin{equation}\label{eq:jsbudget}
D_{\mathrm{JS}}\!\left(
\frac{n\,\bm{p}_{\mathrm{real}} +
\lambda |\tilde{\Dsyn}| \bm{p}_{\tilde{\Dsyn}}}
{n + \lambda|\tilde{\Dsyn}|},\ 
\bm{p}_{\mathrm{real}}\right) \leq \varepsilon,
\end{equation}
其中 $n = |\Dreal|$。预算是\emph{自适应}的：门控按层内标准化置
信度从高到低贪心注入候选，每注入一个样本即重估
\eqref{eq:jsbudget}，超出预算则跳过该样本继续扫描。当候选池健
康（漂移小）时准入更多样本，池恶化时自动收紧。

\subsection{与置信门控的关系}
纯置信门控（仅 \eqref{eq:znorm} 无 \eqref{eq:jsbudget}）控制
\emph{单样本}质量，不控制\emph{集合}分布；纯预算门控反之。二
者互补：前者保证每个注入样本有依据，后者保证注入集合整体不离
真实分布。消融实验（第四章）分别去除两轴验证其必要性。

\section{加权混合与后继训练}\label{sec:mix}
通过门控的样本按权重 $\lambda$ 与真实数据混合。实现上以整数复
制率 $\lceil \lambda \rceil$ 近似（$\lambda \leq 1$ 时取 1，
后续由轮次训练轮数隐式缩放），后继训练从 $M^{(r-1)}$ 参数出发
在 $\Dmix$ 上以较小学习率精调若干轮。权重 $\lambda$ 控制自生成
信号对训练分布的影响强度，其敏感性在第四章分析。

\section{评估反馈}\label{sec:feedback}
每轮在留出测试集上全量排序评估 NDCG@10 与 HR@10，并记录门控统
计（候选数、准入数、混合 JS 距离）。迭代曲线与漂移曲线共同构
成闭环健康监控：若指标随轮次下降或 JS 距离超预算，则提示闭环
失稳。

\section{训练流程与算法}
算法~\ref{alg:rsi} 汇总完整流程。

\begin{algorithm}[t]
\caption{保真度门控 RSI 闭环（单次实验运行）}
\label{alg:rsi}
\begin{algorithmic}[1]
\Require 真实训练集 $\Dreal$，视图数 $K$，阈值 $\tau, m, z_{\min}$，
预算 $\varepsilon$，权重 $\lambda$，轮数 $R$
\State $M^{(0)} \gets \mathrm{Train}(\text{随机初始化}; \Dreal)$
\Comment{预训练至平台}
\State 计算流行度分层 $\sigma(\cdot)$ 与 $\bm{p}_{\mathrm{real}}$
\For{$r = 1, \dots, R$}
  \ForAll{用户 $u$，历史 $S_u$}
    \State 依 \eqref{eq:view} 构造 $K$ 个视图，计算 $\bm{p}^{(k)}$
    \State 若满足 \eqref{eq:consistency}：输出候选 $(u, \hat{i})$
  \EndFor
  \State 对候选做层内 z 标准化 \eqref{eq:znorm}，保留 $z \geq z_{\min}$
  \State 按 $z$ 降序贪心注入，逐样本检查 JS 预算 \eqref{eq:jsbudget}
  \State $\Dmix \gets \Dreal \cup \lambda\tilde{\Dsyn}^{(r)}$
  \State $M^{(r)} \gets \mathrm{Train}(M^{(r-1)}; \Dmix)$
  \Comment{小步长精调}
  \State 评估并记录指标与门控统计
\EndFor
\State \Return 迭代历史 $\{M^{(r)}\}_{r=0}^{R}$ 与全部统计
\end{algorithmic}
\end{algorithm}

\section{复杂度分析}
设用户数 $U$、物品数 $N$、平均序列长度 $\ell$、GRU 前向代价
$O(\ell d^2 + \ell d N)$（$d$ 为隐维度，含输出投影）。每轮：
生成阶段 $K$ 次前向/用户，$O(UK(\ell d^2 + \ell dN))$；一致性
与门控为线性扫描，$O(UKN)$；重训为若干轮全量训练。相对一轮训
练，生成开销约放大 $K$ 倍前向——实验中 $K=3$ 时生成约占每轮总
时间的 $30\%$。空间开销可忽略（候选与统计均为 $O(UK)$）。

\section{误差累积的定性分析}
崩塌的根源是\emph{无界的分布漂移}：每轮漂移 $\delta_r =
D_{\mathrm{JS}}(\bm{p}_{\mathrm{mix}}^{(r)}, \bm{p}_{\mathrm{real}})$
被训练放大并叠加。本文框架从两端施加约束：生成端，一致性条件削
弱头部偏袒，使单轮漂移的\emph{来源}减小；门控端，预算
\eqref{eq:jsbudget} 直接给出逐轮漂移上界 $\delta_r \leq
\varepsilon$。在 $\varepsilon$ 有界且每轮以真实数据为主（
$\lambda |\tilde{\Dsyn}| \ll n$）的设定下，后继训练的分布始终
锚定在 $\Dreal$ 的 $\varepsilon$-邻域内，迭代不再发散。第四章
的漂移曲线与消融结果与此分析一致。

\section{本章小结}
本章给出 RSI 闭环的形式化定义、朴素自训练的崩塌机理（命题
~\ref{prop:collapse}），并提出针对性框架：多视图一致性生成改变
信号的偏袒结构，双轴保真度门控从单样本与集合分布两端控制准入，
JS 预算为逐轮漂移提供显式上界。
